%============================================================
% CHƯƠNG 3: ĐẠO LÀM TRÒ
%============================================================
\chapter{Đạo Làm Trò (Being a Good Student)}

\begin{center}
  \textit{\color{truyenblue}``Education is not the filling of a pail, but the lighting of a fire.''}\\
  \textit{(Giáo dục không phải là đổ đầy một cái thùng, mà là thắp lên một ngọn lửa.)}\\[4pt]
  \small\color{truyengold}--- W.B. Yeats ---
\end{center}
\ngancach

\section{Câu Hỏi Cuối Giờ}
\begin{truyen}{Câu Hỏi Cuối Giờ}{Ham Học Hỏi}
\chuhoa{C}{ }hỉ còn năm phút hết giờ, cả lớp đã đóng vở chuẩn bị ra về. Chỉ có một em giơ tay.\\
\textit{(With only five minutes left, the whole class had already closed their notebooks. Only one student raised a hand.)}

``Thầy ơi, em vẫn chưa hiểu tại sao kết quả lại ra như vậy.''\\
\textit{(``Teacher, I still do not understand why the result came out this way.'')}

Thầy giáo nhìn đồng hồ, nhìn lại em học sinh, rồi mỉm cười ngồi xuống bàn.\\
\textit{(The teacher glanced at the clock, then at the student, then smiled and sat back down.)}

``\textbf{Câu hỏi hay nhất hôm nay. Lại đây thầy giải thích cho}.''\\
\textit{(``\textbf{The best question of today. Come here so I can explain}.'')}

Em học sinh ở lại thêm hai mươi phút. Đó là \textbf{hai mươi phút thay đổi cả một hướng tư duy}.\\
\textit{(The student stayed an extra twenty minutes. Those were \textbf{twenty minutes that changed an entire way of thinking}.)}
\end{truyen}
\begin{baihoc}
  Học trò giỏi không phải là người có nhiều câu trả lời đúng mà là người dám đặt câu hỏi khi chưa hiểu. Tò mò và dũng cảm trí tuệ là chìa khóa mở mọi cánh của tri thức.
\end{baihoc}

\section{Vở Ghi Mượn}
\begin{truyen}{Vở Ghi Mượn}{Trân Trọng}
\chuhoa{E}{ }m học sinh nhà nghèo không có tiền mua vở mới. Người bạn cùng bàn chia đôi cuốn vở của mình.\\
\textit{(The poor student had no money to buy a new notebook. The deskmate split a notebook in half to share.)}

Từng trang ghi chép cẩn thận, nét chữ ngay ngắn như thể cậu \textbf{biết đây là thứ quý giá được nhường cho mình}.\\
\textit{(Every page was noted carefully, the handwriting neat as if he \textbf{knew this was something precious given to him}.)}

Cuối năm học, cậu đứng nhất lớp. Nhận giải, cậu cầm micro và nói:

``\textbf{Tôi muốn cảm ơn người bạn đã chia nửa cuốn vở cho tôi. Đó là nửa cuốn vở nhưng là cả tấm lòng}.''\\
\textit{(``\textbf{I want to thank the friend who shared half a notebook with me. It was half a notebook but a whole heart}.'')}
\end{truyen}
\begin{baihoc}
  Học trò biết ơn và trân trọng cơ hội học tập sẽ không bao giờ lãng phí những gì được trao cho. Biết ơn là sức mạnh đẩy con người tiến xa hơn cả tài năng thiên bẩm.
\end{baihoc}

\section{Bài Kiểm Tra Làm Lại}
\begin{truyen}{Bài Kiểm Tra Làm Lại}{Trung Thực}
\chuhoa{T}{ }rong lúc cô giáo ra ngoài, bạn cùng bàn thì thầm: ``Tao có đáp án rồi. Mày chép đi.''\\
\textit{(While the teacher stepped out, the deskmate whispered: ``I have the answers. Just copy.'')}

Em học sinh nhìn tờ đáp án, nhìn lại bài làm còn bỏ trống của mình, rồi lắc đầu.\\
\textit{(The student looked at the answer sheet, then at her own still-blank test paper, then shook her head.)}

``Tao chưa học kỹ phần này. Điểm thấp để tao biết mình cần học lại.''\\
\textit{(``I did not study this part well. A low grade will tell me what I need to relearn.'')}

Cô giáo hôm sau gọi riêng em ra: ``Em là người duy nhất không có điểm cao bất thường hôm qua. \textbf{Em xứng đáng được thầy cô tin tưởng nhất}.''\\
\textit{(The next day the teacher called her aside: ``You were the only one without an unusually high score yesterday. \textbf{You deserve the most trust from us}.'')}
\end{truyen}
\begin{baihoc}
  Điểm số gian lận chỉ dối người chấm bài, không thể dối chính bản thân. Học trò trung thực với lỗ hổng của mình mới thực sự trưởng thành trong tri thức.
\end{baihoc}

\section{Bài Học Từ Lỗi Sai}
\begin{truyen}{Bài Học Từ Lỗi Sai}{Khiêm Tốn}
\chuhoa{E}{ }m học sinh giỏi nhất lớp trả lời sai một câu trong giờ vật lý. Cả lớp cười.\\
\textit{(The best student in class gave a wrong answer during physics class. The whole class laughed.)}

Em đứng thẳng lưng: ``\textbf{Cảm ơn thầy đã sửa. Giờ em hiểu rồi.}'' Rồi ngồi xuống viết lại vào vở.\\
\textit{(The student stood straight: ``\textbf{Thank you, teacher, for the correction. Now I understand.}'' Then sat down and rewrote it in the notebook.)}

Thầy dừng lại nhìn cả lớp: ``Các em vừa thấy điều gì? Không phải lỗi sai. Mà là \textbf{cách một người học giỏi thực sự xử lý lỗi sai}.''\\
\textit{(The teacher paused and looked at the class: ``What did you all just see? Not a mistake. But \textbf{how a truly good learner handles a mistake}.'')}
\end{truyen}
\begin{baihoc}
  Khiêm tốn trước lỗi lầm không làm giảm đi trí tuệ. Ngược lại, nó chính là biểu hiện của trí tuệ cao nhất – khả năng học hỏi liên tục mà không để tự ái cản đường.
\end{baihoc}

\section{Cuốn Sách Cũ}
\begin{truyen}{Cuốn Sách Cũ}{Kiên Nhẫn}
\chuhoa{C}{ }uốn sách giáo khoa cũ của em bị mất bìa, nhiều trang nhàu nát và có chữ gạch chân khắp nơi.\\
\textit{(The student's old textbook was missing its cover, many pages were crumpled and filled with underlined text throughout.)}

Bạn cùng lớp hỏi: ``Sao mày không mua sách mới?''\\
\textit{(A classmate asked: ``Why not buy a new book?'')}

``\textbf{Cuốn này tao đọc mười lần rồi. Mỗi chỗ gạch chân là một thứ tao đã hiểu ra}. Sách mới thì đẹp nhưng sách này thì thiệt.''\\
\textit{(``\textbf{I have read this one ten times. Every underline is something I came to understand}. A new book looks pretty but this one is real.'')}

Kỳ thi đại học, em đất nhất tỉnh.\\
\textit{(At the university entrance exam, the student scored highest in the province.)}
\end{truyen}
\begin{baihoc}
  Học không cần sách đẹp, cần đầu óc chịu nghiền ngẫm. Kiên nhẫn đọc đi đọc lại một cuốn sách thường tốt hơn đọc qua loa mười cuốn sách.
\end{baihoc}

\section{Tiết Học Buồn Chán}
\begin{truyen}{Tiết Học Buồn Chán}{Tôn Trọng}
\chuhoa{T}{ }iết học địa lý dài và buồn tẻ. Nhiều bạn lén chơi điện thoại. Thầy giáo già biết nhưng không nói.\\
\textit{(The geography lesson was long and dull. Many students secretly played on their phones. The old teacher knew but said nothing.)}

Cuối tiết, thầy dừng lại nhìn cả lớp và nói nhẹ:

``Ba mươi năm dạy học, hôm nay là lần đầu tiên tôi thấy mình không đủ hay để giữ các em ngồi yên. Xin lỗi các em.''\\
\textit{(At the end of the class, the teacher paused and spoke softly: ``Thirty years of teaching, and today is the first time I feel I was not interesting enough to keep you seated. I am sorry.'')}

Cả lớp \textbf{im lặng và xấu hổ}.\\
\textit{(The whole class \textbf{fell silent and felt ashamed}.)}

Từ hôm đó, \textbf{không một ai lấy điện thoại ra trong tiết thầy}.\\
\textit{(From that day on, \textbf{not a single student took out a phone during his class}.)}
\end{truyen}
\begin{baihoc}
  Tôn trọng thầy cô không phải vì quy định mà vì hiểu được công sức và tâm huyết đằng sau mỗi tiết học. Khi ta nhìn thấy con người sau bục giảng, ta sẽ không nỡ làm họ tổn thương.
\end{baihoc}

\section{Nghề Tương Lai}
\begin{truyen}{Nghề Tương Lai}{Định Hướng}
\chuhoa{C}{ }ô giáo hỏi cả lớp: ``Các em muốn làm nghề gì khi lớn lên?''\\
\textit{(The teacher asked the class: ``What do you want to be when you grow up?'')}

Hầu hết trả lời: bác sĩ, kỹ sư, phi công.\\
\textit{(Most answered: doctor, engineer, pilot.)}

Chỉ có một em nói: ``Em muốn làm người tử tế.''\\
\textit{(Only one student said: ``I want to be a decent person.'')}

Cả lớp cười. Cô giáo không cười. Cô gật đầu:

``\textbf{Đó là nghề khó nhất và cũng là nghề quan trọng nhất. Em đã chọn đúng rồi}.''\\
\textit{(The class laughed. The teacher did not laugh. She nodded: ``\textbf{That is the hardest profession and also the most important one. You have chosen rightly}.'')}
\end{truyen}
\begin{baihoc}
  Mục tiêu học tập không chỉ là kiếm được một nghề nghiệp tốt mà còn là trở thành một người có phẩm cách. Hai điều đó không loại trừ nhau, nhưng cái thứ hai quan trọng hơn nhiều.
\end{baihoc}

\section{Tiền Học Phí}
\begin{truyen}{Tiền Học Phí}{Biết Ơn}
\chuhoa{E}{ }m học bổng đầu tiên anh nhận được, anh không mua gì cho bản thân.\\
\textit{(The first scholarship he received, he did not buy anything for himself.)}

Anh gửi về nhà toàn bộ, kèm theo lá thư: ``Ba mẹ ơi, đây là tiền ba mẹ đã đầu tư vào con. \textbf{Con xin hoàn lại một phần nhỏ}.''\\
\textit{(He sent it all home, along with a letter: ``Mom and Dad, this is the money you invested in me. \textbf{I am returning a small part of it}.'')}

Mẹ anh đọc thư, gọi điện khóc: ``Con giữ lấy mà học tiếp. Ba mẹ chưa bao giờ muốn trả lại gì.''\\
\textit{(His mother read the letter, called in tears: ``Keep it for your studies. We never wanted anything back.'')}

``\textbf{Con biết. Nhưng con muốn ba mẹ biết rằng con không học để kiếm tiền cho mình. Con học để xứng đáng với ba mẹ.}''\\
\textit{(``\textbf{I know. But I want you to know that I am not studying to earn money for myself. I am studying to be worthy of you.}'')}
\end{truyen}
\begin{baihoc}
  Người học trò hiếu thảo hiểu rằng mỗi trang sách là một phần công sức của cha mẹ. Học tốt không chỉ vì tương lai bản thân mà còn là cách trả ơn thiêng liêng nhất.
\end{baihoc}

\section{Hai Học Sinh Giỏi}
\begin{truyen}{Hai Học Sinh Giỏi}{Hợp Tác}
\chuhoa{H}{ }ai học sinh giỏi nhất lớp luôn cạnh tranh nhau. Mỗi kỳ thi đều chỉ hơn kém nhau một hai điểm.\\
\textit{(The two best students in class were always competing. Every exam they differed by just one or two points.)}

Một hôm cô giáo ghép họ thành một cặp làm bài thi nhóm.\\
\textit{(One day the teacher paired them together for a group project.)}

Cả hai nghi ngại. Nhưng kết quả: \textbf{đội họ đạt điểm tuyệt đối} – không ai đứng nhất, cả hai cùng vượt qua.\\
\textit{(Both were apprehensive. But the result: \textbf{their team achieved a perfect score} – neither ranked first, yet both surpassed their previous limits.)}

``\textbf{Hóa ra tài năng của mày bù đúng chỗ tao yếu. Và ngược lại}.''\\
\textit{(``\textbf{It turns out your talent fills exactly where I am weak. And vice versa}.'')}
\end{truyen}
\begin{baihoc}
  Cạnh tranh lành mạnh thúc đẩy tiến bộ, nhưng hợp tác tạo ra những kết quả mà cạnh tranh không bao giờ đạt được. Người học trò khôn ngoan biết khi nào cần chiến đấu và khi nào cần hợp sức.
\end{baihoc}

\section{Lời Xin Lỗi Muộn}
\begin{truyen}{Lời Xin Lỗi Muộn}{Sửa Lỗi}
\chuhoa{M}{ }ười năm sau khi ra trường, anh tìm lại người thầy dạy lớp 9 để xin lỗi.\\
\textit{(Ten years after graduation, he went back to find the Grade 9 teacher to apologise.)}

Ngày ấy anh đã cùng bạn bè chép tay những câu chuyện cười về thầy dán lên bảng tin.\\
\textit{(Back then, he and his friends had handwritten jokes about the teacher and pinned them to the notice board.)}

Thầy đã già, tóc bạc, nhưng nhận ra anh ngay.\\
\textit{(The teacher was old now, with white hair, but recognised him immediately.)}

``Thưa thầy, con xin lỗi vì điều con đã làm năm lớp 9.''\\
\textit{(``Teacher, I am sorry for what I did in Grade 9.'')}

Thầy nhìn anh lâu, rồi nói: ``\textbf{Thầy đã tha thứ cho em từ hôm đó. Nhưng cảm ơn em đã dạy thầy rằng học trò cũ vẫn có thể lớn lên}.''\\
\textit{(The teacher looked at him for a long time, then said: ``\textbf{I forgave you on that very day. But thank you for teaching me that former students can still grow up}.'')}
\end{truyen}
\begin{baihoc}
  Xin lỗi muộn còn hơn không xin lỗi. Dũng cảm thừa nhận lỗi lầm quá khứ và tìm cách sửa chữa là dấu hiệu của một nhân cách đã trưởng thành thật sự.
\end{baihoc}
